\documentclass[UTF8,12pt,a4paper]{ctexart}

\usepackage{amsmath}
\usepackage{cases}
\usepackage{cite}
\usepackage{graphicx}
\usepackage{enumerate}
\usepackage{algorithm}
\usepackage{caption} %\caption*需要
\usepackage[noend]{algpseudocode} %algorithmicx
\makeatletter
\renewcommand{\fnum@algorithm}{\fname@algorithm}
\makeatother
\renewcommand{\algorithmicrequire}{\textbf{Input:}}
\renewcommand{\algorithmicensure}{\textbf{Output:}}
\usepackage[margin=1in]{geometry}
\geometry{a4paper}
\usepackage{fancyhdr}
\pagestyle{fancy}
\fancyhf{}
\usepackage{xcolor}
\usepackage{listings}
\lstset{
	basicstyle=\ttfamily,                                % 传统上，代码打字机字体好些
	breaklines,
	tabsize=4,
	extendedchars=false,
	columns=fixed,
	%numbers=left,                                        % 在左侧显示行号
	numbers=none,
	frame=none,                                          % 不显示背景边框
	backgroundcolor=\color[RGB]{245,245,244},            % 设定背景颜色
	keywordstyle=\color[RGB]{40,40,255},                 % 设定关键字颜色
	numberstyle=\footnotesize\color{darkgray},           % 设定行号格式
	commentstyle=\it\color[RGB]{0,96,96},                % 设置代码注释的格式
	stringstyle=\rmfamily\slshape\color[RGB]{128,0,0},   % 设置字符串格式
	showstringspaces=false,                              % 不显示字符串中的空格
	xleftmargin=1em,
	%xrightmargin=1em,
	%aboveskip=-0.5em
	language=c++,                                        % 设置语言
}

\title{算法原理 HW4 作业}
\author{
	姓名：\underline{冯峻}~~~~~~
	学号：\underline{523031910148}~~~~~~}
\date{\today}
\pagenumbering{arabic}

\begin{document}

% 如下两行请按实际情况修改
\fancyhead[L]{冯峻}
\fancyhead[C]{动态规划}
\fancyfoot[C]{\thepage}

\maketitle
\tableofcontents
\newpage

\begin{enumerate}[1.]
	\item \textbf{题目：}给定正整数集合 $S=\{a_1, a_2,\cdots, a_n\}$，设计动态规划算法求解该集合的划分问题，即判断是否存在集合 $A\subseteq S$，使得$\sum_{a_i\in A}a_i=\sum_{a_i\in S-A}a_i$。分析该问题的优化子结构、重叠子问题，给出动态规划算法伪代码并分析算法的复杂度。

	\textbf{解：}

	\textbf{问题分析：}

	该问题等价于判断是否存在子集$A$，使得其元素之和等于集合总和的一半。设集合$S$的总和为$sum$，若$sum$为奇数，则显然无解；若$sum$为偶数，问题转化为：是否存在子集，其元素之和等于$target = sum/2$。这是经典的\textbf{0-1背包问题}的判定版本。

	\textbf{优化子结构：}

	定义子问题：$dp[i][j]$表示使用前$i$个元素能否凑出和为$j$。对于第$i$个元素$a_i$，有两种选择：
	\begin{itemize}
		\item 不选$a_i$：问题规模缩减为前$i-1$个元素能否凑出和为$j$，即$dp[i][j] = dp[i-1][j]$
		\item 选$a_i$：问题规模缩减为前$i-1$个元素能否凑出和为$j-a_i$，即$dp[i][j] = dp[i-1][j-a_i]$
	\end{itemize}
	
	因此，原问题的最优解包含子问题的最优解，具有\textbf{优化子结构}性质。

	\textbf{重叠子问题：}

	在计算$dp[i][j]$时，可能需要多次计算相同的子问题$dp[i-1][j']$。例如，计算$dp[3][10]$和$dp[3][12]$时，都可能需要$dp[2][8]$的结果。这种重复计算的子问题体现了\textbf{重叠子问题}性质。

	\textbf{动态规划算法伪代码：}
	\begin{algorithm}[H]%
		\caption{SUBSET-PARTITION($S$)}
		\begin{algorithmic}[1]
			\State $sum \gets \sum_{i=1}^{n} S[i]$
			\If{$sum \bmod 2 \neq 0$}
				\State \Return{False}
			\EndIf
			\State $target \gets sum / 2$
			\State 初始化二维数组$dp[0..n][0..target]$，全部置为False
			\State $dp[i][0] \gets$ True, $\forall i \in [0, n]$ \Comment{和为0总是可以达到}
			\For{$i = 1$ to $n$}
				\For{$j = 1$ to $target$}
					\State $dp[i][j] \gets dp[i-1][j]$ \Comment{不选第$i$个元素}
					\If{$j \geq S[i]$}
						\State $dp[i][j] \gets dp[i][j]$ OR $dp[i-1][j-S[i]]$ \Comment{选第$i$个元素}
					\EndIf
				\EndFor
			\EndFor
			\State \Return{$dp[n][target]$}
		\end{algorithmic}
	\end{algorithm}

	\textbf{时间复杂度分析：}

	算法需要填充一个$n \times target$的二维表，其中$target = sum/2$。每个状态的计算时间为$O(1)$。因此，总时间复杂度为$O(n \cdot target) = O(n \cdot sum)$，其中$sum$为集合总和。

	\textbf{空间复杂度分析：}

	使用二维数组的空间复杂度为$O(n \cdot target)$。可以优化为一维数组（滚动数组），空间复杂度降为$O(target) = O(sum)$。

	\textbf{C++代码实现见：}code/subset\_partition.cpp

	\newpage

	\item \textbf{题目：}课上讲过了子串和子序列的区别。现考虑两个字符串$X$和$Y$，求解其最长的公共子串。比如$X=$'photograph'，$Y=$'tomography'，则其最长公共子串为'ograph'。
	
	1) 已知$n=|X|$，$m=|Y|$，请设计一个时间复杂度$\Theta(nm)$的动态规划算法求解$X$和$Y$的最长公共子串。

	2) 设计一个时间复杂度$\Theta(nm)$的非动态规划算法求解$X$和$Y$的最长公共子串。

	需给出相应的算法设计和分析过程，不能仅仅只给出伪代码。

	\textbf{解：}

	\textbf{问题理解：}

	首先明确\textbf{子串}和\textbf{子序列}的区别：
	\begin{itemize}
		\item \textbf{子串}：必须是连续的字符序列
		\item \textbf{子序列}：可以不连续，但必须保持相对顺序
	\end{itemize}

	本题求解的是\textbf{最长公共子串}（Longest Common Substring），要求字符必须连续出现。

	\subsection*{1) 动态规划算法}

	\textbf{优化子结构分析：}

	定义状态$dp[i][j]$表示\textbf{以$X[i-1]$和$Y[j-1]$结尾}的最长公共子串的长度。关键是"以...结尾"这个限定条件。

	状态转移方程：
	$$dp[i][j] = \begin{cases}
		dp[i-1][j-1] + 1, & \text{if } X[i-1] = Y[j-1] \\
		0, & \text{if } X[i-1] \neq Y[j-1]
	\end{cases}$$

	当$X[i-1] = Y[j-1]$时，以它们结尾的最长公共子串长度等于以$X[i-2]$和$Y[j-2]$结尾的最长公共子串长度加1。当字符不相等时，由于子串必须连续，长度为0。

	这与最长公共\textbf{子序列}不同。在子序列问题中，当字符不相等时，状态转移为$dp[i][j] = \max(dp[i-1][j], dp[i][j-1])$；而在子串问题中，由于连续性要求，状态直接置为0。

	\textbf{重叠子问题：}

	在计算$dp[i][j]$时，需要用到$dp[i-1][j-1]$的结果。多个状态的计算会重复使用相同的子问题结果，体现了重叠子问题性质。

	\textbf{动态规划算法伪代码：}
	\begin{algorithm}[H]%
		\caption{LCS-SUBSTRING-DP($X, Y$)}
		\begin{algorithmic}[1]
			\State $n \gets |X|$, $m \gets |Y|$
			\State 初始化二维数组$dp[0..n][0..m]$，全部置为0
			\State $maxLength \gets 0$, $endIndex \gets 0$
			\For{$i = 1$ to $n$}
				\For{$j = 1$ to $m$}
					\If{$X[i-1] = Y[j-1]$}
						\State $dp[i][j] \gets dp[i-1][j-1] + 1$
						\If{$dp[i][j] > maxLength$}
							\State $maxLength \gets dp[i][j]$
							\State $endIndex \gets i-1$ \Comment{记录在$X$中的结束位置}
						\EndIf
					\Else
						\State $dp[i][j] \gets 0$
					\EndIf
				\EndFor
			\EndFor
			\State 从$X[endIndex - maxLength + 1..endIndex]$提取最长公共子串
			\State \Return{最长公共子串及其长度$maxLength$}
		\end{algorithmic}
	\end{algorithm}

	\textbf{时间复杂度：}$\Theta(nm)$，需要填充$n \times m$的二维表，每个状态计算为$O(1)$。

	\textbf{空间复杂度：}$O(nm)$，可优化至$O(m)$（使用滚动数组）。

	\textbf{C++代码实现见：}code/longest\_common\_substring\_dp.cpp

	\subsection*{2) 非动态规划算法}

	\textbf{算法设计思路：}

	采用\textbf{滑动窗口法}（Two Pointers），直接在两个字符串上进行匹配比较。

	\textbf{核心思想：}
	\begin{itemize}
		\item 枚举$X$中的每个起始位置$i$
		\item 枚举$Y$中的每个起始位置$j$
		\item 从位置$(i, j)$开始，尽可能向后匹配更长的公共子串
		\item 记录所有匹配过程中的最长子串
	\end{itemize}

	\textbf{算法正确性：}

	该算法遍历了所有可能的子串起始位置组合，对于任意最长公共子串，必然存在某个$(i, j)$组合，使得从该位置开始的匹配能够覆盖到它。因此算法能够正确找到最长公共子串。

	\textbf{非动态规划算法伪代码：}
	\begin{algorithm}[H]%
		\caption{LCS-SUBSTRING-SLIDING($X, Y$)}
		\begin{algorithmic}[1]
			\State $n \gets |X|$, $m \gets |Y|$
			\State $maxLength \gets 0$, $endIndex \gets 0$
			\For{$i = 0$ to $n-1$} \Comment{枚举$X$的起始位置}
				\For{$j = 0$ to $m-1$} \Comment{枚举$Y$的起始位置}
					\State $length \gets 0$
					\State $x \gets i$, $y \gets j$
					\While{$x < n$ AND $y < m$ AND $X[x] = Y[y]$}
						\State $length \gets length + 1$
						\State $x \gets x + 1$, $y \gets y + 1$
					\EndWhile
					\If{$length > maxLength$}
						\State $maxLength \gets length$
						\State $endIndex \gets i + length - 1$
					\EndIf
				\EndFor
			\EndFor
			\State 从$X[endIndex - maxLength + 1..endIndex]$提取最长公共子串
			\State \Return{最长公共子串及其长度$maxLength$}
		\end{algorithmic}
	\end{algorithm}

	\textbf{时间复杂度分析：}

	外层两个循环分别遍历$X$和$Y$的所有起始位置，时间复杂度为$O(nm)$。内层while循环最多执行$\min(n, m)$次。

	在最坏情况下（如两个完全相同的字符串），时间复杂度为$O(nm \cdot \min(n, m))$，这似乎超过了$\Theta(nm)$的要求。

	但我们注意到，所有内层while循环的总执行次数实际上是有界的。每次while循环会在字符串上"滑动"，对于任意一对字符$(X[x], Y[y])$，它最多被比较常数次。更精确的分析表明，总的字符比较次数为$O(nm)$。

	因此，该算法的时间复杂度为$\Theta(nm)$。

	\textbf{空间复杂度：}$O(1)$，只使用常数个变量。

	\textbf{C++代码实现见：}code/longest\_common\_substring\_non\_dp.cpp

\end{enumerate}


\end{document}
